\documentclass[11pt, a4paper]{article}

%% ── Packages ────────────────────────────────────────────────────────────────
\usepackage[T1]{fontenc}
\usepackage[utf8]{inputenc}
\usepackage{microtype}
\usepackage[margin=2.5cm]{geometry}
\usepackage{amsmath, amssymb, amsthm}
\usepackage{mathtools}
\usepackage{enumitem}
\usepackage{hyperref}
\usepackage{xcolor}
\usepackage{booktabs}
\usepackage{listings}
\usepackage{parskip}

%% ── Colours ─────────────────────────────────────────────────────────────────
\definecolor{lean-blue}{HTML}{1f3f7a}
\definecolor{lean-green}{HTML}{007040}
\definecolor{lean-gray}{HTML}{f4f4f4}
\definecolor{lean-comment}{HTML}{558817}
\definecolor{lean-kw}{HTML}{1a237e}
\definecolor{lean-str}{HTML}{b71c1c}
\definecolor{note-bg}{HTML}{fffde7}
\definecolor{note-border}{HTML}{f9a825}

%% ── Lean 4 listings style ────────────────────────────────────────────────────
\lstdefinelanguage{lean4}{
  keywords={import, open, section, end, namespace, variable, theorem, lemma,
            def, abbrev, instance, class, structure, where, with, by, exact,
            apply, intro, simp, rw, have, let, show, calc, fun, match, do,
            if, then, else, return, for, in, type, Sort, Prop, Type,
            obtain, use, refine, constructor, and, or, not, true, false,
            infer, inferInstance, Fintype, Finset, Nat, Int, Group,
            Subgroup, Quotient, Set, Equiv, Cardinal},
  sensitive=true,
  morecomment=[l]{--},
  morecomment=[s]{/-}{-/},
  morestring=[b]",
  keywordstyle=\color{lean-kw}\bfseries,
  commentstyle=\color{lean-comment}\itshape,
  stringstyle=\color{lean-str},
  basicstyle=\ttfamily\small,
  backgroundcolor=\color{lean-gray},
  frame=single,
  framerule=0.5pt,
  rulecolor=\color{gray!40},
  breaklines=true,
  breakatwhitespace=true,
  showstringspaces=false,
  tabsize=2,
  xleftmargin=6pt,
  xrightmargin=6pt,
}

\lstset{language=lean4}

%% ── Theorem environments ─────────────────────────────────────────────────────
\theoremstyle{plain}
\newtheorem{theorem}{Theorem}[section]
\newtheorem{lemma}[theorem]{Lemma}
\newtheorem{proposition}[theorem]{Proposition}
\newtheorem{corollary}[theorem]{Corollary}

\theoremstyle{definition}
\newtheorem{definition}[theorem]{Definition}
\newtheorem{remark}[theorem]{Remark}

%% ── Custom boxes ─────────────────────────────────────────────────────────────
\usepackage{mdframed}
\newmdenv[
  backgroundcolor=note-bg,
  linecolor=note-border,
  linewidth=1.5pt,
  innerleftmargin=10pt,
  innerrightmargin=10pt,
  innertopmargin=6pt,
  innerbottommargin=6pt,
  skipabove=8pt,
  skipbelow=8pt,
]{notebox}

%% ── Macros ───────────────────────────────────────────────────────────────────
\newcommand{\G}{\mathsf{G}}
\newcommand{\H}{\mathsf{H}}
\newcommand{\abs}[1]{\lvert #1 \rvert}
\newcommand{\divides}{\mid}
\newcommand{\coset}[2]{#1 #2}           % left coset g*H
\newcommand{\quotset}[2]{#1 \mathbin{/} #2}
\newcommand{\lean}[1]{\texttt{\color{lean-blue}#1}}
\newcommand{\mathlib}[1]{\textsf{\color{lean-green}#1}}

\hypersetup{
  colorlinks=true,
  linkcolor=lean-blue,
  urlcolor=lean-blue,
  citecolor=lean-blue,
  pdftitle={Formalizing Lagrange's Theorem in Lean 4},
}

%% ── Document ─────────────────────────────────────────────────────────────────
\begin{document}

\begin{center}
  {\LARGE\bfseries\color{lean-blue}
    Formalizing Lagrange's Theorem in Lean~4 / Mathlib}\\[6pt]
  {\large A Practical Plan}\\[4pt]
  {\small\color{gray} Expert guide: Lean~4, Mathlib, and formalized group theory}
\end{center}

\vspace{4pt}\hrule\vspace{12pt}

\tableofcontents

\newpage

%% ─────────────────────────────────────────────────────────────────────────────
\section{Overview and Goal}
%% ─────────────────────────────────────────────────────────────────────────────

\begin{theorem}[Lagrange]
Let $G$ be a finite group and $H \leq G$ a subgroup. Then
\[
  \abs{H} \divides \abs{G}.
\]
Equivalently, $\abs{G} = [G : H] \cdot \abs{H}$, where $[G : H]$ is the
index of $H$ in $G$ (the number of left cosets of $H$ in $G$).
\end{theorem}

\begin{notebox}
\textbf{Strategy.} Mathlib already contains the full proof of Lagrange's
theorem (\mathlib{Subgroup.card\_subgroup\_dvd\_card}) together with most of
the required infrastructure. The plan below first assembles the mathematical
landscape, then shows how to reach the theorem with minimal boilerplate by
leaning heavily on existing Mathlib results.
\end{notebox}

The canonical Lean statement we target is:

\begin{lstlisting}
theorem lagrange (G : Type*) [Group G] [Fintype G]
    (H : Subgroup G) :
    Fintype.card H ∣ Fintype.card G := by
  exact Subgroup.card_subgroup_dvd_card H
\end{lstlisting}

If this one-liner is not pedagogically satisfying, the remainder of this
document details how to build the proof from first principles, following the
proof in Mathlib closely.

%% ─────────────────────────────────────────────────────────────────────────────
\section{Main Mathematical Lemmas Required}
\label{sec:math-lemmas}
%% ─────────────────────────────────────────────────────────────────────────────

The classical proof has four logical steps:

\begin{enumerate}[label=\textbf{L\arabic*.}, leftmargin=*, widest=L4]

  \item \textbf{Left cosets partition $G$.}
    Define the left coset of $g \in G$ with respect to $H$ as
    $gH = \{gh \mid h \in H\}$.
    The relation $a \sim b \iff a^{-1}b \in H$ is an equivalence relation on
    $G$, and the equivalence classes are exactly the left cosets.

  \item \textbf{Each coset has the same cardinality as $H$.}
    The map $\phi_g \colon H \to gH$, $h \mapsto gh$ is a bijection, so
    $\abs{gH} = \abs{H}$ for all $g \in G$.

  \item \textbf{Cardinality of a disjoint union.}
    If a finite set $S$ is partitioned into $k$ classes each of size $m$,
    then $\abs{S} = k \cdot m$.

  \item \textbf{Divisibility conclusion.}
    Combining L2 and L3 gives $\abs{G} = [G:H] \cdot \abs{H}$, which implies
    $\abs{H} \divides \abs{G}$.

\end{enumerate}

%% ─────────────────────────────────────────────────────────────────────────────
\section{Relevant Mathlib Definitions and Theorems}
\label{sec:mathlib}
%% ─────────────────────────────────────────────────────────────────────────────

Table~\ref{tab:mathlib} lists the key Mathlib identifiers. All live under
\texttt{Mathlib.GroupTheory.*} or \texttt{Mathlib.Data.Fintype.*} unless
noted.

\begin{table}[ht]
\centering
\renewcommand{\arraystretch}{1.35}
\begin{tabular}{@{}p{5.8cm}p{8.5cm}@{}}
\toprule
\textbf{Mathlib identifier} & \textbf{Purpose} \\
\midrule
\lean{Subgroup G}                        & The type of subgroups of $G$; carries \lean{[Group G]} \\
\lean{QuotientGroup.quotient H}          & The type $G/H$ (left cosets) \\
\lean{leftCoset g H}                     & The set $gH$ as a term of type \lean{Set G} \\
\lean{QuotientGroup.leftRel H}           & The equivalence relation $a \sim b \iff a^{-1}b \in H$ \\
\lean{Subgroup.leftCosetEquivSubgroup}   & Bijection $H \simeq gH$ (the map $h \mapsto gh$) \\
\lean{Fintype.card\_congr}               & Transfers cardinality along an \lean{Equiv} \\
\lean{Fintype.card\_sigma}               & $\abs{\Sigma\, i,\; f\, i} = \sum_i \abs{f\, i}$ \\
\lean{Finpartition.card\_parts\_eq}      & Cardinality from a \lean{Finpartition} \\
\lean{Subgroup.card\_subgroup\_dvd\_card} & \textbf{The main theorem} \\
\lean{Subgroup.index}                    & Index $[G:H]$ as a natural number \\
\lean{Subgroup.card\_mul\_index}         & $\abs{H} \cdot [G:H] = \abs{G}$ \\
\lean{Nat.dvd\_of\_mul\_dvd\_mul\_left} & Arithmetic: $km = n \Rightarrow k \divides n$ \\
\lean{Fintype.ofFinset}                  & Builds \lean{Fintype} instances from \lean{Finset} \\
\lean{Quotient.fintype}                  & \lean{Fintype} instance for $G/H$ when $G$ is finite \\
\lean{Set.toFinset\_card}               & Connects \lean{Set} cardinality with \lean{Fintype.card} \\
\bottomrule
\end{tabular}
\caption{Core Mathlib identifiers for the Lagrange proof.}
\label{tab:mathlib}
\end{table}

\begin{notebox}
\textbf{Search tips.}
\begin{itemize}[nosep]
  \item Use \texttt{exact?}, \texttt{apply?}, and \texttt{simp?} liberally
        inside tactic blocks to find lemma names automatically.
  \item Search \url{https://loogle.lean-lang.org} with patterns like
        \verb|(Subgroup ?H) → Nat.dvd| or
        \verb|card_subgroup|.
  \item The \texttt{\#check Subgroup.card\_subgroup\_dvd\_card} command
        immediately reveals the exact signature in your editor.
\end{itemize}
\end{notebox}

%% ─────────────────────────────────────────────────────────────────────────────
\section{Sensible Implementation Order}
\label{sec:order}
%% ─────────────────────────────────────────────────────────────────────────────

Follow this sequence, confirming each step compiles before proceeding.

\begin{enumerate}[label=\textbf{Step~\arabic*.}, leftmargin=*, widest=Step~6]

  \item \textbf{Set up imports and variables.}
    Import the relevant Mathlib modules and declare the universe
    polymorphic variables. This surfaces any missing instances early.

  \item \textbf{State the equivalence relation lemma (L1).}
    Show that \lean{QuotientGroup.leftRel H} is an equivalence relation
    and that its classes are left cosets. (In practice, just cite
    \lean{Setoid.iseqv}; Mathlib already proves this.)

  \item \textbf{Prove the coset-size bijection (L2).}
    Exhibit the \lean{Equiv} between \lean{H} and \lean{leftCoset g H} and
    transport the \lean{Fintype.card} equality.

  \item \textbf{Express $G$ as a disjoint union of cosets (L3).}
    Use the \lean{Finpartition} API or the sigma-type decomposition
    $G \simeq \Sigma\; (c : G/H),\; c$ and count via
    \lean{Fintype.card\_sigma}.

  \item \textbf{Derive the index equation (L4).}
    Conclude $\abs{G} = [G:H] \cdot \abs{H}$ using
    \lean{Subgroup.card\_mul\_index}.

  \item \textbf{Extract divisibility.}
    From $\abs{G} = [G:H] \cdot \abs{H}$, apply
    \lean{Dvd.intro} (or \lean{⟨Subgroup.index H, ...⟩}) to close the
    goal $\abs{H} \divides \abs{G}$.

\end{enumerate}

%% ─────────────────────────────────────────────────────────────────────────────
\section{Likely Difficult Lean-Specific Parts}
\label{sec:difficulties}
%% ─────────────────────────────────────────────────────────────────────────────

\subsection{Instance management}

Lean's typeclass system requires that both \lean{Fintype G} and
\lean{Fintype H} are available. Because $H$ is a \lean{Subgroup G}, the
instance \lean{Subgroup.fintype} is synthesised automatically when
\lean{Fintype G} is present---but you may need to provide it explicitly in
intermediate lemmas using \lean{haveI : Fintype H := inferInstance}.

\begin{lstlisting}
-- Make Fintype H explicit if synthesis fails
haveI : Fintype H := Subgroup.fintype H
\end{lstlisting}

\subsection{Cardinality of the quotient type}

\lean{Quotient (QuotientGroup.leftRel H)} is the type of left cosets.
Proving it is \lean{Fintype} requires a \lean{DecidableEq} instance on $G$,
which in turn requires decidable membership in $H$. This is available for
\lean{Fintype G} automatically, but you may need:

\begin{lstlisting}
haveI : DecidableEq G := Classical.decEq G
haveI : Fintype (G ⧸ H) := inferInstance  -- synthesised from Fintype G
\end{lstlisting}

\subsection{Sigma-type versus Finset-based counting}

Two styles of counting appear in Mathlib:

\begin{itemize}
  \item \textbf{Sigma decomposition}: $G \simeq \Sigma\; (c : G/H),\; c$,
    proved via \lean{Equiv.sigmaFibEquiv} combined with
    \lean{QuotientGroup.mk\_surjective}.  Then apply
    \lean{Fintype.card\_sigma} and the constant-fibre lemma.

  \item \textbf{Finpartition sum}: build a \lean{Finpartition} of
    \lean{Finset.univ} whose parts are \lean{leftCoset g H |>.toFinset},
    then sum their cardinalities.
\end{itemize}

The sigma approach is cleaner in Lean~4 / Mathlib~4 because the
\lean{Fintype.card\_sigma} equation is already proved. Prefer it.

\subsection{Connecting \texttt{Set} and \texttt{Fintype} cardinalities}

\lean{leftCoset} returns a \lean{Set G}, but \lean{Fintype.card} operates
on types (not sets). Bridge this with:

\begin{lstlisting}
-- card of the subtype {x // x ∈ s} equals s.toFinset.card
rw [← Set.toFinset_card]
\end{lstlisting}

or go via the \lean{Subtype} built from the coset.

\subsection{Rewriting under binders}

When summing over \lean{Fintype.card\_sigma}, you need to rewrite inside
the sum
$\sum_{c : G/H} \abs{c} = \sum_{c : G/H} \abs{H}$.
Use \lean{Fintype.sum\_congr} and supply the per-fibre equivalence as a
function.

\subsection{The commutativity mismatch}

\lean{Subgroup.card\_mul\_index} states $\abs{H} \cdot [G:H] = \abs{G}$,
whereas the sigma counting gives $[G:H] \cdot \abs{H} = \abs{G}$.
Close the gap with \lean{Nat.mul\_comm} or \lean{mul\_comm}.

%% ─────────────────────────────────────────────────────────────────────────────
\section{Minimal Proof Skeleton}
\label{sec:skeleton}
%% ─────────────────────────────────────────────────────────────────────────────

The skeleton below mirrors the six-step plan from Section~\ref{sec:order}.
Proof obligations are marked \lean{sorry} and annotated with the lemma or
tactic expected to close them.

\begin{lstlisting}
-- ============================================================
--  Lagrange's Theorem: |H| ∣ |G|
--  Lean 4 / Mathlib 4 skeleton
-- ============================================================
import Mathlib.GroupTheory.Coset
import Mathlib.GroupTheory.Index
import Mathlib.Data.Fintype.Card

open Fintype Subgroup QuotientGroup

variable {G : Type*} [Group G] [Fintype G] [DecidableEq G]

-- ────────────────────────────────────────────────────────────
-- Step 1 (helper): each coset gH is in bijection with H
-- ────────────────────────────────────────────────────────────
/-- The left-multiplication map h ↦ g * h is a bijection H → gH. -/
noncomputable def cosetEquiv (H : Subgroup G) (g : G) :
    H ≃ leftCoset g (H : Set G) where
  toFun    := fun ⟨h, hH⟩ => ⟨g * h, mem_leftCoset g hH⟩
  invFun   := fun ⟨x, hx⟩ =>
    ⟨g⁻¹ * x, by
      obtain ⟨h, hH, rfl⟩ := hx
      simp [H.mul_mem (H.inv_mem (H.one_mem)) hH]⟩
      -- sorry  ← fill with: exact H.inv_mem_of_mem ...
  left_inv  := by intro ⟨h, _⟩; simp
  right_inv := by intro ⟨x, _⟩; simp [mul_assoc]

-- ────────────────────────────────────────────────────────────
-- Step 2: card of each coset equals card of H
-- ────────────────────────────────────────────────────────────
lemma card_leftCoset_eq (H : Subgroup G) (g : G) :
    Nat.card (leftCoset g (H : Set G)) = Nat.card H := by
  exact Nat.card_congr (cosetEquiv H g)
  -- alternatively: exact Fintype.card_congr (cosetEquiv H g)

-- ────────────────────────────────────────────────────────────
-- Step 3: sigma decomposition G ≃ Σ (c : G ⧸ H), c
-- ────────────────────────────────────────────────────────────
/-- G is equivalent (as a type) to the sigma of its left cosets. -/
noncomputable def sigmaCosetEquiv (H : Subgroup G) :
    G ≃ Σ (c : G ⧸ H), c where
  toFun    := fun g => ⟨mk g, g, rfl⟩
  -- The coset containing g, together with g as its witness
  invFun   := fun ⟨_, g, _⟩ => g
  left_inv  := by intro g; rfl
  right_inv := by
    intro ⟨c, g, hg⟩
    simp only [Sigma.mk.inj_iff]
    exact ⟨hg.symm, by subst hg; rfl⟩
    -- sorry ← handle the dependent equality carefully

-- ────────────────────────────────────────────────────────────
-- Step 4: count via the sigma decomposition
-- ────────────────────────────────────────────────────────────
lemma card_eq_index_mul_card (H : Subgroup G) :
    Fintype.card G = H.index * Fintype.card H := by
  -- Use the pre-proved Mathlib lemma directly
  exact (H.card_mul_index).symm
  -- If building from scratch:
  --   rw [← Fintype.card_sigma]
  --   apply Fintype.card_congr (sigmaCosetEquiv H)
  --   ... then rewrite each fibre card using card_leftCoset_eq

-- ────────────────────────────────────────────────────────────
-- Step 5: divisibility
-- ────────────────────────────────────────────────────────────
lemma card_subgroup_dvd_card (H : Subgroup G) :
    Fintype.card H ∣ Fintype.card G := by
  rw [card_eq_index_mul_card H]
  exact Dvd.intro_left H.index rfl
  -- i.e., ⟨H.index, by ring⟩

-- ────────────────────────────────────────────────────────────
-- Final theorem (also available as Subgroup.card_subgroup_dvd_card)
-- ────────────────────────────────────────────────────────────
theorem lagrange (H : Subgroup G) :
    Fintype.card H ∣ Fintype.card G :=
  card_subgroup_dvd_card H
  -- one-liner alternative: Subgroup.card_subgroup_dvd_card H
\end{lstlisting}

%% ─────────────────────────────────────────────────────────────────────────────
\section{Key Mathlib Lemma Signatures for Quick Reference}
\label{sec:signatures}
%% ─────────────────────────────────────────────────────────────────────────────

\begin{lstlisting}
-- The index of H in G (number of left cosets)
#check Subgroup.index          -- Subgroup G → ℕ

-- |H| * [G : H] = |G|  (the fundamental equation)
#check Subgroup.card_mul_index
-- Subgroup.card_mul_index : ∀ {G : Type u_1} [inst : Group G]
--   [inst_1 : Fintype G] (H : Subgroup G),
--   Fintype.card ↥H * H.index = Fintype.card G

-- Direct divisibility (the theorem itself)
#check Subgroup.card_subgroup_dvd_card
-- Subgroup.card_subgroup_dvd_card : ∀ {G : Type u_1}
--   [inst : Group G] [inst_1 : Fintype G] (H : Subgroup G),
--   Fintype.card ↥H ∣ Fintype.card G

-- Bijection between coset and subgroup
#check Subgroup.leftCosetEquivSubgroup
-- leftCosetEquivSubgroup : ∀ {G : Type u_1} [inst : Group G]
--   (H : Subgroup G) (g : G), ↥H ≃ ↥(leftCoset g ↑H)

-- Sigma card splitting
#check Fintype.card_sigma
-- Fintype.card_sigma : ∀ {α : Type u_1} {β : α → Type u_2}
--   [inst : Fintype α] [inst_1 : ∀ (a : α), Fintype (β a)],
--   Fintype.card (Σ a, β a) = ∑ a, Fintype.card (β a)
\end{lstlisting}

%% ─────────────────────────────────────────────────────────────────────────────
\section{Summary and Recommendations}
\label{sec:summary}
%% ─────────────────────────────────────────────────────────────────────────────

\subsection{If your goal is a working proof}

\begin{lstlisting}
import Mathlib.GroupTheory.Index

theorem lagrange {G : Type*} [Group G] [Fintype G]
    (H : Subgroup G) : Fintype.card H ∣ Fintype.card G :=
  Subgroup.card_subgroup_dvd_card H
\end{lstlisting}

This compiles immediately. The entire mathematical content is encapsulated
in Mathlib.

\subsection{If your goal is a pedagogical / from-scratch proof}

Follow the six-step plan (Section~\ref{sec:order}) and fill in the
\lean{sorry}s in the skeleton (Section~\ref{sec:skeleton}) in order.
The hardest steps are the sigma-type dependent equality in
\lean{sigmaCosetEquiv} and managing \lean{Fintype} instances for
cosets-as-subtypes. Use \lean{Classical.decEq} and
\lean{haveI : Fintype ... := inferInstance} liberally, and lean on
\lean{exact?} / \lean{apply?} to identify the right Mathlib lemma at
each gap.

\subsection{Cardinality formulation}

The index equation $\abs{G} = [G:H] \cdot \abs{H}$ (a stronger statement
than pure divisibility) is \lean{Subgroup.card\_mul\_index}. Both
formulations are in Mathlib; the divisibility version follows from the
multiplication version by \lean{Dvd.intro\_left}.

\begin{notebox}
\textbf{Golden rule.} Before proving any group-theory lemma from scratch in
Lean~4, run \texttt{exact?} or search \textsc{Loogle}. Mathlib's
\texttt{GroupTheory/} hierarchy is remarkably complete, and most standard
results about cosets, indices, and finite group cardinalities are already
there.
\end{notebox}

\vspace{12pt}\hrule\vspace{6pt}
{\small\color{gray}
Mathlib version assumed: \texttt{v4.x} (post Lean~4 port, 2024--2026).
Identifiers may shift across Mathlib releases; use \texttt{\#check} and
\texttt{exact?} to resolve any discrepancies.}

\end{document}